\chapter{Konzeption des Assistenzsystems}
\label{chap:konzeption}

Das vorliegende Kapitel beschreibt die konzeptionelle Grundlage des entwickelten Assistenzsystems. 
Ausgehend von den in Kapitel~4 ermittelten Anforderungen werden zunächst die Zielarchitektur 
und die Datenpipeline erläutert, bevor das Modell der Kontextualisierung und Priorisierung sowie 
das gewählte Ausgabeformat vorgestellt werden. Jede Entwurfsentscheidung wird dabei mit 
Bezug auf die identifizierten funktionalen und nicht-funktionalen Anforderungen begründet.

% ─────────────────────────────────────────────────────────────────────────────
\section{Zielarchitektur und Datenpipeline}
\label{sec:zielarchitektur}

Die Zielarchitektur des Systems folgt dem Prinzip einer sequenziellen Datenpipeline, 
die drei voneinander unabhängige Analysequellen zusammenführt und deren Ergebnisse 
in einem einheitlichen Schema normalisiert, bevor sie an ein lokales Sprachmodell 
übergeben werden. Abbildung~\ref{fig:pipeline} illustriert den Gesamtablauf.

% TODO: Abbildung der Pipeline hier einfügen
% \begin{figure}[h]
%   \centering
%   \includegraphics[width=\textwidth]{figures/pipeline.pdf}
%   \caption{Datenpipeline des ACE-Assistenzsystems}
%   \label{fig:pipeline}
% \end{figure}

Die Pipeline gliedert sich in vier Phasen: die parallele Datenerhebung durch axe-core, 
Playwright und grep, die Normalisierung der Rohdaten in das Unified Finding Schema (UFS), 
die Prompt-Konstruktion mit Token-Budget-Steuerung sowie die Inferenz durch das lokale 
Sprachmodell und die anschließende Ausgabe.

\subsection{Analysequellen}

Das System nutzt drei komplementäre Werkzeuge zur Datenerhebung, die jeweils 
unterschiedliche Aspekte der Barrierefreiheit abdecken und sich gegenseitig ergänzen.

\textbf{axe-core} wird über einen gesteuerten Chromium-Browser ausgeführt und analysiert 
den vollständig gerenderten DOM der Zielanwendung auf Verstöße gegen WCAG-Erfolgskriterien. 
Da Rich-Client-Webanwendungen ihre Inhalte erst zur Laufzeit per JavaScript erzeugen, 
ist ein echter Browser-Launch unerlässlich -- statische HTML-Analysen würden den 
tatsächlichen Zustand der Anwendung nicht abbilden.\footcite[Vgl.][]{fathallah_accessguru_2025}

\textbf{Playwright} führt sieben generische Interaktionschecks aus, die unabhängig von 
der konkreten Anwendungslogik auf jeder Web-App lauffähig sind. Geprüft werden unter 
anderem die Erreichbarkeit aller interaktiven Elemente per Tastatur, die Sichtbarkeit 
des Fokusindikators sowie das Vorhandensein eines Skip-Links. Diese verhaltensbasierten 
Checks decken Verletzungen auf, die im statischen DOM nicht sichtbar sind, etwa fehlende 
Fokus-Traps in modalen Dialogen.\footcite[Vgl.][]{huang_access_2024}

\textbf{grep} durchsucht den React-Quellcode der Anwendung nach bekannten Anti-Patterns 
und reichert gleichzeitig axe-core-Findings um ihren Entstehungskontext im Quellcode an. 
Die Entscheidung für reguläre Ausdrücke anstelle eines vollständigen Abstract Syntax Tree 
(AST)-Parsers ergibt sich aus zwei Überlegungen: Erstens sind die sechs implementierten 
Pattern-Checks syntaktischer Natur -- sie suchen nach textuell erkennbaren Mustern im 
JSX-Markup, nicht nach semantischen Datenflüssen, für die ein AST-Parser erforderlich wäre. 
Zweitens gewährleistet der regexp-basierte Ansatz Sprachunabhängigkeit über 
\texttt{.tsx}, \texttt{.jsx}, \texttt{.ts} und \texttt{.js}-Dateien hinweg, ohne 
sprachspezifische Parser-Abhängigkeiten einzuführen.\footcite[Vgl.][S.~12]{bassi_2025}
Die Anreicherungsstrategie folgt dabei einem selektiven Ansatz: Durchsucht werden 
ausschließlich jene Quelldateien, auf die axe-core oder Playwright hinweisen, um das 
verfügbare Token-Budget nicht durch irrelevante Codefragmente zu belasten (vgl. 
Abschnitt~\ref{sec:kontextualisierung}).

% ─────────────────────────────────────────────────────────────────────────────
\subsection{Unified Finding Schema}
\label{sec:ufs}

Ein zentrales architektonisches Element ist das \emph{Unified Finding Schema} (UFS), 
das alle Findings unabhängig von ihrer Herkunft in eine einheitliche TypeScript-Schnittstelle 
überführt. Ohne eine solche Normalisierungsschicht würde der Prompt-Builder drei 
strukturell heterogene Datenformate erhalten: axe-core liefert Objekte mit 
\texttt{impact}-Stufen und DOM-Selektorketten, Playwright gibt boolesche 
\texttt{passed}-Felder mit Check-IDs zurück, und grep produziert Datei-Zeilen-Tupel 
mit Textfragmenten. Diese Heterogenität hätte drei negative Konsequenzen.

Erstens würde die Token-Effizienz sinken, da rohe axe-core-Ausgaben vollständige 
DOM-Snapshots und ausführliche Hilfetexte enthalten, die für die LLM-Inferenz nicht 
relevant sind. Zweitens wäre eine quellenübergreifende Deduplizierung nicht möglich, 
da kein gemeinsames Identifikationsfeld existierte. Drittens könnte die 
Severity-basierte Priorisierung im Token-Budget-Algorithmus nicht angewendet werden, 
da \texttt{impact} (axe), \texttt{severity} (intern) und Dateiname (grep) 
inkommensurable Größen darstellen.\footcite[Vgl.][]{schulhoff_prompt_2024}

Das UFS definiert folgende Kernfelder:

\begin{itemize}
  \item \texttt{source}: Herkunftsquelle (\texttt{"axe"}, \texttt{"playwright"}, \texttt{"grep"})
  \item \texttt{ruleId}: Eindeutiger Regelbezeichner, normalisiert über alle Quellen
  \item \texttt{severity}: Einheitliche Schweregradbewertung (\texttt{critical}, \texttt{serious}, \texttt{moderate}, \texttt{minor})
  \item \texttt{wcagCriteria} und \texttt{wcagLevel}: Zugeordnetes WCAG-Erfolgskriterium und Konformitätsstufe
  \item \texttt{selector}: CSS-Selektor des betroffenen DOM-Elements
  \item \texttt{componentPath}: Dateipfad der React-Komponente (sofern durch grep ermittelt)
  \item \texttt{codeSnippet} und \texttt{lineRange}: Relevanter Quellcodeausschnitt mit Zeilenangabe
  \item \texttt{category}: Klassifikation nach \texttt{"syntaktisch"}, \texttt{"semantisch"} oder \texttt{"layout"} in Anlehnung an die Taxonomie von \footcite[Vgl.][]{fathallah_accessguru_2025}
\end{itemize}

Das Feld \texttt{componentPath} ist dabei das entscheidende Bindeglied zwischen 
DOM-Analyse und Quellcode: Erst durch die Verknüpfung eines axe-Findings mit dem 
konkreten Dateipfad und der Zeilennummer wird aus einer generischen Meldung 
(\emph{,,Element hat keinen alternativen Text``}) eine entwicklernahe Handlungsanweisung 
(\emph{,,In \texttt{LoginForm.tsx}, Zeile~46: Füge \texttt{alt}-Attribut hinzu``}).

% ─────────────────────────────────────────────────────────────────────────────
\section{Modell der Kontextualisierung und Priorisierung}
\label{sec:kontextualisierung}

\subsection{Single-Prompt-Strategie}

Eine grundlegende Entwurfsentscheidung betrifft die Interaktionsstrategie mit dem 
lokalen Sprachmodell. Konzeptionell wäre ein mehrstufiger Ansatz denkbar, bei dem 
zunächst die axe-Findings, dann die Playwright-Ergebnisse und abschließend der 
Quellcodekontext in separaten Aufrufen verarbeitet werden, wie es etwa 
\citeauthor{fathallah_accessguru_2025} mit dem AccessGuru-System über eine 
Cloud-API realisieren.\footcite[Vgl.][S.~5]{fathallah_accessguru_2025}

Für den vorliegenden Anwendungsfall wird hingegen eine Single-Prompt-Strategie gewählt, 
die alle drei Datenquellen in einem einzigen Modellaufruf zusammenführt. Diese 
Entscheidung ergibt sich aus drei Überlegungen.

\textbf{Latenzbudget:} Die Laufzeitmessungen zeigen, dass ein einzelner LLM-Aufruf 
auf einem lokalen CPU-basierten System zwischen 491 und 3.238 Sekunden benötigt, 
abhängig von der Modellgröße. Eine Multi-Step-Kette mit vier sequenziellen Aufrufen 
würde die Gesamtlaufzeit auf ein Vielfaches erhöhen und damit NFA-02 
(maximale Analysedauer von 20 Minuten) noch deutlicher verletzen, als es beim 
größten getesteten Modell ohnehin der Fall ist.

\textbf{Cross-Source-Korrelation:} Ein einzelner Prompt ermöglicht dem Sprachmodell, 
Zusammenhänge zwischen den Quellen zu erkennen und zu einem kohärenten Fix 
zusammenzufassen. Meldet axe-core ein fehlendes \texttt{aria-label} und liefert grep 
gleichzeitig den zugehörigen Quellcodeausschnitt, kann das Modell beides in einem 
einzigen, konkreten Vorschlag verbinden. Bei getrennten Aufrufen ginge dieser 
Kontext verloren.\footcite[Vgl.][]{schulhoff_prompt_2024}

\textbf{Auditierbarkeit:} Ein deterministischer, reproduzierbarer Prompt ist für eine 
wissenschaftliche Evaluation wichtiger als hypothetische Qualitätsgewinne durch 
mehrstufige Verarbeitung. Ein Prompt entspricht einem nachvollziehbaren 
Entscheidungspfad; Fehler in einem Teilschritt propagieren sich bei Multi-Step 
unkontrolliert in nachfolgende Aufrufe.\footcite[Vgl.][]{njifenjou_role_2024}

\subsection{Token-Budget-Steuerung}

Da lokale Modelle typischerweise über ein Kontextfenster von 32.768 Tokens verfügen 
und der System-Prompt sowie das Ausgabeformat einen fixen Anteil belegen, steht für 
die Findings ein variables Input-Budget zur Verfügung. Der Prompt-Builder implementiert 
eine Priorisierungsstrategie, die Findings nach Schweregrad sortiert und Einträge 
niedrigerer Priorität auslässt, sobald das Budget überschritten wird. Die Sortierung 
erfolgt auf Basis des normierten \texttt{severity}-Felds des UFS, was die 
quellenübergreifende Vergleichbarkeit der Schweregrade voraussetzt.

Darüber hinaus wird grep selektiv eingesetzt: Statt die gesamte Codebasis zu 
durchsuchen, werden ausschließlich Dateien analysiert, die durch axe-core- oder 
Playwright-Findings referenziert werden. Dieses Vorgehen begrenzt die Anzahl 
der Code-Snippets im Prompt und verhindert eine unkontrollierte Ausweitung des 
Token-Verbrauchs.

\subsection{Modellwahl}

Als Inferenzbackend wird Ollama eingesetzt, das die lokale Ausführung von 
Quantisierungsmodellen ohne GPU-Beschleunigung ermöglicht und damit die 
Datenschutzanforderungen der BITBW erfüllt (NFA-05).\footcite[Vgl.][]{paternò_2025}
Cloud-basierte Dienste wie die OpenAI API scheiden aus, da sie eine Übermittlung 
potenziell schützenswerter Quellcodeausschnitte an externe Server erforderten, 
was im Kontext öffentlicher Verwaltung nicht zulässig ist.

Für den Modellvergleich in der Evaluation (vgl. Kapitel~7) wurden drei Modelle 
ausgewählt: \texttt{qwen2.5-coder:7b} als ressourcenschonende Baseline mit 
Fokus auf Code-Aufgaben, \texttt{qwen3:latest} (ca.~8B Parameter) als 
Vertreter der neueren Modellgeneration mit Reasoning-Modus sowie 
\texttt{qwen3:32b} als qualitativ stärkste lokale Option innerhalb des 
verfügbaren Arbeitsspeichers von 96~GB. Die Modellgröße wird dabei bewusst 
als Variable im Evaluationsdesign behandelt, um den Trade-off zwischen 
Ausgabequalität und Laufzeit empirisch zu erfassen.

% ─────────────────────────────────────────────────────────────────────────────
\section{Ausgabeformat: entwicklernahe To-Do-Liste}
\label{sec:ausgabeformat}

Das Ausgabeformat des Systems ist als priorisierte To-Do-Liste konzipiert, 
die sich an Entwicklerinnen und Entwickler richtet und unmittelbar in den 
bestehenden Entwicklungsworkflow integrierbar sein soll. Die Wahl dieses Formats 
ergibt sich aus der Anforderungsanalyse (FA-05) sowie aus der Literatur, die 
zeigt, dass die Akzeptanz von Barrierefreiheitsempfehlungen stark von ihrer 
Handlungsorientierung abhängt.\footcite[Vgl.][S.~8]{fathallah_accessguru_2025}

Die Liste gliedert sich in zwei Prioritätsstufen: \emph{Must-have}-Einträge 
bezeichnen Verstöße gegen Level-A- und Level-AA-Erfolgskriterien, die eine 
gesetzliche Konformitätspflicht begründen können, während \emph{Nice-to-have}-Einträge 
Empfehlungen zur Verbesserung der Nutzungserfahrung umfassen. Jeder Eintrag 
enthält -- sofern durch die grep-Anreicherung verfügbar -- den Dateipfad der 
betroffenen Komponente, die Zeilennummer sowie einen konkreten Code-Fix-Vorschlag. 
Dieser Aufbau entspricht dem Qualitätskriterium der \emph{Kontextsensitivität}: 
Das System geht über die reine Fehlerbenennung hinaus und liefert den 
entwicklungsrelevanten Kontext, der für eine unmittelbare Behebung erforderlich ist.

Es ist anzumerken, dass das Ausgabeformat durch einen System-Prompt vorgegeben 
wird, dessen Einhaltung durch das Sprachmodell nicht garantiert werden kann. 
Die Evaluation in Kapitel~7 zeigt, dass alle drei getesteten Modelle das 
vorgegebene Must-have/Nice-to-have-Schema in mindestens einem Lauf nicht 
einhalten. Diese Limitation wird in Abschnitt~\ref{sec:limitationen} diskutiert.